\documentclass[11pt]{article}
\usepackage[T1]{fontenc}
\usepackage{geometry}
\geometry{margin=1.15in}
\usepackage{amsmath,amssymb,amsthm}
\newtheorem{prop}{Proposition}
\DeclareMathOperator{\had}{h}
\newcommand{\eps}{\varepsilon}
\title{The partition-size hypothesis in Theorem 1.5}
\author{}
\date{}
\begin{document}
\maketitle

All logarithms are natural. We write $\had(G)$ for the largest integer $t$ such that $G$ has a $K_t$ minor, and $\kappa(G)$ for the vertex connectivity of $G$.

\begin{prop}
Let $\eps=e^{-32}$. For every sufficiently large even integer $n$, there is a graph $G$ on $n$ vertices and a partition $(X_1,\dots,X_k)$ of $V(G)$ into sets of size two such that $\kappa(G)\ge\eps n$ and
\[
 \had(G)<
 \frac{n}{\sqrt{\log n}}
 \left(
 \sqrt{-\sum_{i,j=1}^k\alpha_i\alpha_j\log(1-p_{ij}^*)}
 -\eps
 \right),
\]
where $\alpha_i=|X_i|/n$, $p_{ij}$ is the density of edges between $X_i$ and $X_j$,
\[
 p_{ii}=\frac{e(G[X_i])}{\binom{|X_i|}{2}},
 \qquad
 p_{ij}^*=\min\{p_{ij},1-\eps\}.
\]
In particular, the conclusion of Theorem 1.5 fails for unrestricted partitions.
\end{prop}

\begin{proof}
Let $n$ be even, put $k=n/2$, and fix a partition $(X_1,\dots,X_k)$ of $[n]$ into sets of size two. Let $G$ be the random graph on $[n]$ in which each edge is present independently with probability $1/2$. We will show that, with probability tending to one as $n\to\infty$ through even integers,
\begin{equation}\label{eq:three-properties}
 \kappa(G)\ge\frac n{10},
 \qquad
 \had(G)<\frac n{\sqrt{\log n}},
 \qquad
 \Lambda^*:=-\sum_{i,j=1}^k\alpha_i\alpha_j\log(1-p_{ij}^*)\ge\frac32.
\end{equation}

First consider connectivity. For any two distinct vertices $u,v$, the number $Y_{uv}$ of their common neighbors has distribution $\operatorname{Bin}(n-2,1/4)$, with mean $\mu=(n-2)/4$. For $n\ge10$, we have $n/10\le\mu/2$. The Chernoff bound therefore gives
\[
 \Pr(Y_{uv}<n/10)
 \le \Pr(Y_{uv}\le\mu/2)
 \le e^{-\mu/8}
 =e^{-(n-2)/32}.
\]
A union bound over all pairs shows that, with probability tending to one, every two distinct vertices have at least $n/10$ common neighbors. On this event, after deleting any set of fewer than $n/10$ vertices, every two remaining vertices still have a common neighbor. The remaining graph is connected, proving the first assertion in \eqref{eq:three-properties}.

Next put
\[
 s=\left\lceil\frac n{\sqrt{\log n}}\right\rceil.
\]
An ordered partition $(A_1,\dots,A_s)$ of $[n]$ into nonempty sets is called \emph{permissible} for $G$ if, for every $1\le a<b\le s$, at least one edge of $G$ has one endpoint in $A_a$ and the other in $A_b$.

If $\had(G)\ge s$, choose $s$ disjoint connected branch sets of a $K_s$ minor and distribute all remaining vertices among them. Adding vertices preserves the existence of an edge between every two sets, so this gives a permissible ordered partition of $[n]$.

Fix an ordered partition $(A_1,\dots,A_s)$, and put $a_i=|A_i|$ and $M=\binom{s}{2}$. The events that an edge exists between $A_i$ and $A_j$ involve disjoint sets of random edges for different unordered pairs $\{i,j\}$, and hence are independent. Thus
\begin{align*}
 \Pr((A_1,\dots,A_s)\text{ is permissible})
 &=\prod_{1\le i<j\le s}(1-2^{-a_i a_j})\\
 &\le\exp\left(-\sum_{1\le i<j\le s}2^{-a_i a_j}\right).
\end{align*}
Since
\[
 \sum_{i<j}a_i a_j
 =\frac12\left(n^2-\sum_i a_i^2\right)
 \le\frac{n^2}{2},
\]
the convexity of $x\mapsto2^{-x}$ gives
\[
 \sum_{i<j}2^{-a_i a_j}
 \ge M\,2^{-\frac1M\sum_{i<j}a_i a_j}
 \ge M\,2^{-n^2/[s(s-1)]}.
\]
Every ordered partition of $[n]$ into $s$ nonempty parts is represented by a surjection $[n]\to[s]$. Since there are $s^n$ maps $[n]\to[s]$ altogether, there are at most $s^n$ such ordered partitions. The union bound now gives
\begin{equation}\label{eq:minor-bound}
 \Pr(\had(G)\ge s)
 \le
 \exp\left(
 n\log s-\binom{s}{2}
 \exp\left(-\frac{(\log 2)n^2}{s(s-1)}\right)
 \right).
\end{equation}
By the choice of $s$,
\[
 \frac{n^2}{s(s-1)}=(1+o(1))\log n,
 \qquad
 \binom{s}{2}=(1+o(1))\frac{n^2}{2\log n}.
\]
Consequently the subtracted term in the exponent of \eqref{eq:minor-bound} is
\[
 \frac{n^{\,2-\log 2+o(1)}}{2\log n}.
\]
Its ratio to $n\log n$ tends to infinity, because $1-\log 2>0$. Since $n\log s\le n\log n$ for all sufficiently large $n$, the right side of \eqref{eq:minor-bound} tends to zero. Hence, with probability tending to one,
\[
 \had(G)\le s-1<\frac n{\sqrt{\log n}},
\]
as required.

Finally, let $Z$ be the number of unordered pairs $\{i,j\}$ of distinct classes for which all four possible edges between $X_i$ and $X_j$ are present. These events are independent for different unordered pairs, each with probability $1/16$, so
\[
 Z\sim\operatorname{Bin}\left(\binom{k}{2},\frac1{16}\right),
 \qquad
 \mathbb E Z=\frac{k(k-1)}{32},
 \qquad
 \operatorname{Var}(Z)\le\mathbb E Z.
\]
It follows from Chebyshev's inequality that $Z/k^2$ converges in probability to $1/32$. For each pair counted by $Z$, we have $p_{ij}=p_{ji}=1$, and therefore
\[
 -\log(1-p_{ij}^*)=-\log\eps=32.
\]
Here $\alpha_i=1/k$ for every $i$. The sum defining $\Lambda^*$ is over ordered pairs, so each pair counted by $Z$ contributes $64/k^2$. Every other summand is nonnegative. Hence
\[
 \Lambda^*\ge\frac{64Z}{k^2}.
\]
The right side converges in probability to $2$, and thus $\Pr(\Lambda^*\ge3/2)\to1$. This proves the last assertion in \eqref{eq:three-properties}.

All three assertions in \eqref{eq:three-properties} therefore hold simultaneously with probability tending to one. In particular, for every sufficiently large even $n$, there is a graph satisfying them. Since
\[
 \eps=e^{-32}<\frac1{10}
 \qquad\text{and}\qquad
 \sqrt{\frac32}-\eps>1,
\]
this graph has $\kappa(G)\ge\eps n$ but
\[
 \had(G)
 <\frac n{\sqrt{\log n}}
 <\frac n{\sqrt{\log n}}\left(\sqrt{\Lambda^*}-\eps\right).
\]
This proves the proposition.
\end{proof}
\end{document}
